\documentclass[11pt,a4paper]{article}

% -----------------------------------------------------------------------------
% Paquetes: todos forman parte de una instalacion estandar de LaTeX.
% -----------------------------------------------------------------------------
\usepackage[utf8]{inputenc}
\usepackage{enumitem}
\usepackage{amsmath}
\usepackage[T1]{fontenc}
\usepackage[spanish]{babel}
% Fecha en espanol sin depender de babel ni de spanish.ldf.
\newcommand{\FechaEnEspanol}{%
  \number\day\space de\space
  \ifcase\month
    \or enero\or febrero\or marzo\or abril\or mayo\or junio\or
    julio\or agosto\or septiembre\or octubre\or noviembre\or diciembre
  \fi\space de\space\number\year
}
\usepackage[a4paper,margin=2.5cm,headheight=42pt]{geometry}
\usepackage{graphicx}
\usepackage{xcolor}
\usepackage{fancyhdr}

% lastpage no es necesario para compilar, pero permite mostrar "pagina de N".
\newif\ifHayUltimaPagina
\IfFileExists{lastpage.sty}{
  \usepackage{lastpage}
  \HayUltimaPaginatrue
}{
  \HayUltimaPaginafalse
}

% -----------------------------------------------------------------------------
% CONFIGURACION DEL TRABAJO: editar solamente este bloque para reutilizarlo.
% -----------------------------------------------------------------------------
\newcommand{\NombreMateria}{Redes de Computadoras}
\newcommand{\TipoTrabajo}{Trabajo Práctico}
\newcommand{\NumeroTrabajo}{2}
% Cambiar por "Teórico" o "Práctico" según corresponda.
\newcommand{\TipoTP}{Teórico}
\newcommand{\NombreTrabajo}{Medios de transmisión}
\newcommand{\NombreGrupo}{NetRunners}
\newcommand{\Docente}{FACUNDO NICOLAS OLIVA CUNEO}
\newcommand{\LogoArchivo}{logo.png}
\newcommand{\LogoRuta}{\LogoArchivo}

% Una linea por integrante. Se pueden agregar o quitar filas.
\newcommand{\Integrantes}{%
  \begin{tabular}{ll}
    Bastida, Lucas Ramiro & DNI 39444511 \\
    Contessi, Ruben Andres & DNI 33315235 \\
    Garay, Ignacio Hernan & DNI 44191925 \\
    Giraudo, Juan Pablo & DNI 34970089 \\
    Peñaloza, Gonzalo Adrian & DNI 40441355 \\
    Quinteros del Castillo, Tomás Agustín & DNI 46169739 \\
    Vasconsellos Blason, Benjamin Alejandro & DNI 45642879 \\
  \end{tabular}%
}

% -----------------------------------------------------------------------------
% Elementos comunes de caratula y encabezado.
% -----------------------------------------------------------------------------
\newcommand{\InsertarLogo}[1][0.24\textwidth]{%
  \ifx\LogoRuta\empty
    \fbox{\parbox[c][1cm][c]{#1}{\centering Logo institucional\\\small (colocar archivo)}}%
  \else
    \includegraphics[width=#1,keepaspectratio]{\LogoRuta}%
  \fi
}

\newcommand{\TituloCorto}{TP \NumeroTrabajo (\TipoTP): \NombreTrabajo}

\pagestyle{fancy}
\fancyhf{}
\fancyhead[L]{\InsertarLogo[1.4cm]}
\fancyhead[C]{\small\textbf{\TituloCorto}}
\fancyhead[R]{\small\NombreGrupo}
\renewcommand{\headrulewidth}{0.4pt}
\fancyfoot[C]{\small Pagina \thepage\ifHayUltimaPagina\ de \pageref{LastPage}\fi}
\renewcommand{\footrulewidth}{0pt}

\begin{document}

\begin{titlepage}
  \thispagestyle{empty}
  \begin{center}
    \InsertarLogo[0.32\textwidth]

    \vfill
    {\Large\textsc{\NombreMateria}\par}
    \vspace{1.5cm}
    {\Huge\bfseries \TipoTrabajo\par}
    \vspace{0.35cm}
    {\LARGE N.\textsuperscript{o} \NumeroTrabajo (\TipoTP)\par}
    \vspace{0.8cm}
    {\Large\NombreTrabajo\par}

    \vfill
    \begin{tabular}{rl}
      \textbf{Grupo:}   & \NombreGrupo \\
      \textbf{Docente:} & \Docente
    \end{tabular}

    \vspace{1cm}
    \begin{center}
      \textbf{Integrantes}\\[0.2cm]
      \Integrantes
    \end{center}

    \vfill
    {\large \FechaEnEspanol\par}
  \end{center}
\end{titlepage}

\section*{Cuestiones de repaso}

\begin{enumerate}[label=4.\arabic*), itemsep=0.3cm]

  \item
        \textbf{¿Por qué hay dos cables en un par trenzado de cobre?} \\
        Un enlace de conexión que se realiza con un par trenzado usa dos conductores ya que al momento de conectar los dispositivos a sus puntas estos formarán un bucle cerrado en el que la señal viajará como variaciones de corriente en ese bucle. Los dispositivos sensarán esas variaciones de corriente, que representarán la codificación de la señal que se envía.
  \item
        \textbf{¿Cuáles son las limitaciones del par trenzado?} \\
        Las principales limitaciones del par trenzado son su menor distancia de transmisión, menor ancho de banda y menor velocidad respecto de otros medios guiados. Además, es muy susceptible a interferencias y ruido, debido a su acoplamiento con campos electromagnéticos externos.
  \item
        \textbf{¿Cuál es la diferencia entre el par trenzado no apantallado y el par trenzado apantallado?} \\
        La diferencia principal radica en que el par trenzado apantallado consta de una malla interna en la que está embutido, la cual evita la acción del ruido provocado por ondas electromagnéticas; pero la contrapartida es que es más costoso y difícil de ensamblar.
  \item
        \textbf {Describir los principales componentes del cable de fibra óptica.}\\
        Un cable de fibra óptica posee una estructura cilíndrica concéntrica compuesta por tres elementos principales:
        \begin{itemize}
          \item \textbf{Núcleo (\textit{Core}):} Es la parte central del cable, extremadamente delgada, constituida por una sección cilíndrica de vidrio o plástico con un alto índice de refracción. Por aquí es por donde viajan las ondas de luz conducidas internamente.

          \item \textbf{Revestimiento (\textit{Cladding}):} Es la capa intermedia que rodea al núcleo. Está hecha de un material de vidrio o plástico con propiedades ópticas diferentes. Esta diferencia de índices de refracción atrapa la luz dentro del núcleo mediante el fenómeno de reflexión interna total.

          \item \textbf{Cubierta / Revestimiento exterior (\textit{Jacket}):} Es la capa más externa, fabricada generalmente de materiales plásticos o sintéticos. Su función no es óptica, sino puramente mecánica y de protección física: resguarda al núcleo y al revestimiento contra la humedad, abrasiones, tensiones mecánicas y agresiones del entorno.
        \end{itemize}
  \item \textbf{¿Qué ventajas y desventajas tiene la transmisión de microondas?}

        \textbf{Ventajas:}
        \begin{itemize}
          \item Las microondas permiten realizar enlaces de larga distancia sin necesidad de un medio físico.
          \item En los sistemas de microondas, los amplificadores o repetidores pueden estar más separados entre sí, generalmente entre \(10\) km y \(100\) km.
          \item Las frecuencias más altas permiten utilizar antenas más pequeñas y más baratas, por lo que son adecuadas para distancias más cortas.
          \item Las microondas también pueden utilizarse en comunicaciones vía satélite, permitiendo establecer enlaces entre estaciones terrestres alejadas.
        \end{itemize}

        \textbf{Desventajas:}
        \begin{itemize}
          \item La atenuación aumenta con la lluvia, siendo especialmente significativa para frecuencias superiores a \(10\) GHz.
          \item Debido a la creciente popularidad de las microondas, las áreas de cobertura pueden solaparse, haciendo que las interferencias sean un peligro potencial.
          \item Las bandas de \(4\) a \(6\) GHz, utilizadas habitualmente en transmisiones de larga distancia, sufren una creciente congestión.
          \item Las bandas de frecuencias superiores son menos útiles para distancias largas debido a que la atenuación aumenta.
        \end{itemize}
  \item \textbf{¿Qué es la difusión directa por satélite (DBS, Direct Broadcast Satellite)?} \\
        La difusión directa por satélite (DBS, Direct Broadcast Satellite) es una tecnología de distribución de televisión en la que la señal de vídeo se transmite directamente desde el satélite hasta los domicilios de los usuarios finales.
  \item
        \textbf{¿Por qué un satélite debe usar frecuencias ascendentes y descendentes distintas?} \\
        Debe usar frecuencias distintas porque, en una transmisión continua y sin interferencias, el satélite no puede transmitir y recibir en el mismo rango de frecuencias. Se produciría una interferencia entre la señal entrante y la saliente.
  \item
        \textbf{Indique las diferencias más significativas entre la difusión de radio y las microondas.}
        Diferencias principales: la difusión de microondas es parabólica, es decir direccional, mientras que las ondas de radio son omnidireccionales; como consecuencia de lo anterior, la difusión por microondas requiere de antenas parabólicas bien direccionadas, en cambio las ondas de radio no; como otra diferencia clave está el hecho de que la difusión por microondas es útil para los enlaces punto a punto (antena a antena), en cambio en el caso de la difusión por ondas de radio el uso práctico son los enlaces de un emisor a muchos receptores; también podemos decir que los fenómenos físicos que afectan estas comunicaciones van a ser distintos, por ejemplo las ondas de radio son menos sensibles a la atenuación causada por lluvia que las microondas pero en contrapartida se ven más afectadas por multitrayectorias.
  \item
        \textbf{¿Qué dos funciones realiza una antena?}

        Una antena cumple dos funciones principales según si está transmitiendo o recibiendo la señal:

        \begin{itemize}
          \item \textbf{En emisión:} Recibe la corriente eléctrica del equipo emisor y la transforma en ondas electromagnéticas que se propagan por el aire.
          \item \textbf{En recepción:} Captura las ondas electromagnéticas que viajan por el aire y las transforma nuevamente en una corriente eléctrica para que el equipo receptor pueda procesarla.
        \end{itemize}
  \item
        \textbf{¿Qué es una antena isotrópica?}

        Una antena isotrópica es un punto en el espacio que radia potencia de igual forma en todas las direcciones. Su diagrama de radiación consiste en una esfera centrada en la posición de la antena isotrópica.
  \item
        \textbf{¿Cuál es la ventaja de una antena parabólica por reflexión?} \\
        La ventaja es el direccionamiento de la señal en la emisión y en la recepción. Debido a la forma de la antena una señal emitida desde el foco de la parábola y reflejarse en la antena tendrá una dirección de transmisión paralela al eje de la antena, esto generará un haz de señal que podrá viajar casi en línea recta hacia donde apunte la misma. Además, al recibir una señal con haces que lleguen paralelos al eje de la antena, al reflejarse en la parábola todas incidirán en el foco de la parábola, lo que permite concentrar en un solo punto mayor potencia de la señal.
  \item \textbf{¿Qué factores determinan la ganancia de una antena?} \\
        Los factores físicos y matemáticos que determinan la ganancia de una antena son principalmente el área efectiva y la frecuencia de operación (o longitud de onda).

        \begin{itemize}
          \item \textbf{Área efectiva $(A_e)$:} está relacionada con el tamaño físico, la geometría y la eficiencia de la antena. En una antena parabólica, un mayor diámetro permite obtener una mayor directividad y, en consecuencia, una mayor ganancia. El área efectiva se puede expresar como:

                \[
                  A_e = \eta_a A
                \]

                donde $\eta_a$ es la eficiencia de apertura.

          \item \textbf{Frecuencia de operación $(f)$ o longitud de onda $(\lambda)$:} para un área efectiva constante, la ganancia es inversamente proporcional al cuadrado de la longitud de onda y directamente proporcional al cuadrado de la frecuencia. Por lo tanto, para una misma antena física, una mayor frecuencia de operación generalmente permite obtener una mayor ganancia.
        \end{itemize}

        La relación fundamental es:

        \[
          \boxed{
            G = \frac{4\pi A_e}{\lambda^2}
            = \frac{4\pi f^2 A_e}{c^2}
          }
        \]

        donde $c$ es la velocidad de propagación de las ondas electromagnéticas en el vacío, aproximadamente:

        \[
          c \approx 3 \times 10^8 \, \text{m/s}
        \]
  \item
        \textbf{¿Cuál es la principal causa de la pérdida de señal en comunicaciones vía satélite?} \\
        La principal causa de pérdida de señal en las comunicaciones vía satélite es la pérdida en el espacio libre, producida por la dispersión de la señal debido a la gran distancia entre las antenas.
  \item
        \textbf{¿Qué es la refracción?} \\
        Cuando una onda electromagnética pasa de un medio con una densidad a otro con densidad distinta, su velocidad cambia. El efecto de esto es desviar la dirección de la onda en la separación entre los dos medios. Al pasar de un medio menos denso a otro con densidad mayor, la onda se desviará hacia el medio más denso. Este fenómeno se puede observar sumergiendo parcialmente un palo en agua. La refracción se produce debido a que la velocidad de las ondas electromagnéticas es una función de la densidad del medio atravesado.

        Las ondas de radio se pueden refractar (o desviar) cuando se propagan a través de la atmósfera. La refracción es causada por los cambios en la velocidad de la señal al cambiar su altura o por otro tipo de cambios espaciales debido a las condiciones atmosféricas. Normalmente, la velocidad de la señal aumenta con la altura, haciendo que las ondas de radio se desvíen hacia la superficie terrestre. las condiciones metereológicas pueden implicar variaciones
        en la velocidad con la altura que sean significativamente distintas de las variaciones típicas esperadas. Esto puede dar lugar a que sólo una fracción, o incluso nada, de la onda transmitida siguiendo la trayectoria visual llegue a la antena receptora.
  \item
        \textbf{¿Qué diferencia hay entre difracción y dispersión?} \\
        Para diferenciar, decimos que la Dispersión es el cambio de dirección o frecuencia que sufre una onda al encontrarse con partículas de materia (como gotas de lluvia, por ejemplo), y esto hace que la energía se esparza en muchas direcciones; en cambio, la difracción es un fenómeno que provoca que la onda se curve alrededor del borde de un objeto grande (dada la diferencia de tamaño relativo del objeto y la longitud de onda).
\end{enumerate}

\section*{Ejercicios}


\begin{enumerate}[label=4.\arabic*), itemsep=0.3cm]
  \item
        \textbf{Supóngase que unos datos se almacenan en disquetes de
        \(1,4\) Mbytes que pesan \(30\) g cada uno y que una compañía aérea
        transporta \(10^4\) kg de disquetes a una velocidad de \(1000\) km/h
        sobre una distancia de \(5000\) km. ¿Cuál es la velocidad de transmisión
        en bits por segundo de este sistema?}

        Primero calculamos la cantidad de disquetes transportados:

        \[
        10^4\text{ kg}=10^7\text{ g}
        \]

        \[
        N=\frac{10^7}{30}=333333,33\text{ disquetes}
        \]

        La cantidad de información transportada es:

        \[
        D=333333,33(1,4\times10^6)\text{ bytes}
        \]

        Pasando a bits:

        \[
        D=333333,33(1,4\times10^6)(8)
        \]

        \[
        D=3,7333\times10^{12}\text{ bits}
        \]

        El tiempo necesario para recorrer los \(5000\) km es:

        \[
        t=\frac{5000}{1000}=5\text{ h}
        \]

        \[
        t=5(3600)=18000\text{ s}
        \]

        Por lo tanto, la velocidad de transmisión equivalente es:

        \[
        R=\frac{3,7333\times10^{12}}{18000}
        \]

        \[
        \boxed{R\approx2,07\times10^8\text{ bit/s}}
        \]

        Por lo tanto, la velocidad de transmisión es aproximadamente:

        \[
        \boxed{207,4\text{ Mbit/s}}
        \]

  \item
        Para calcular, primero debemos conocer la potencia de salida.

        \[
          \text{Pérdida} = 10 \log_{10}\left(\frac{P_{\text{in}}}{P_{\text{out}}}\right)
        \]

        \[
          \text{Pérdida} = 20; \quad P_{\text{in}} = 0,5\ \text{W}
        \]

        \[
          20 = 10 \log_{10}\left(\frac{0,5\ \text{W}}{P_{\text{out}}}\right)
        \]

        \[
          \frac{20}{10} = \log_{10}\left(\frac{0,5\ \text{W}}{P_{\text{out}}}\right)
        \]

        \[
          10^{2} = \frac{0,5\ \text{W}}{P_{\text{out}}}
        \]

        \[
          P_{\text{out}} = \frac{0,5\ \text{W}}{100} = 5\ \text{mW}
        \]

        Ahora podemos calcular SNR:

        \[
          \text{SNR} = 10 \log_{10}\left(\frac{P_{\text{señ}}}{P_{\text{rui}}}\right)
        \]

        \[
          P_{\text{señ}} = P_{\text{out}} = 5\ \text{mW}; \quad P_{\text{rui}} = 4,5\ \mu\text{W}
        \]

        \[
          \text{SNR} = 10 \log_{10}\left(\frac{5\ \text{mW}}{4,5\ \mu\text{W}}\right)
        \]

        \[
          \text{SNR} = 10 \log_{10}(1111,11)
        \]

        \[
          \text{SNR} = 10 \times 3,0457
        \]

        \[
          \text{SNR} = 30,457
        \]

        \[
          \boxed{\text{SNR} \approx 30\ \text{dB}}
        \]

  \item
  Para resolver este ejercicio, debemos tener presente dos cosas: la primera, la fórmula de la que vamos a partir es la del cálculo de presupuesto de pérdida:
\[
  L_{dB} = 10\log_{10}\left(\frac{P_{entrada}}{P_{salida}}\right)
\]
Y como segunda cosa, tenemos que ver el gráfico de las curvas de atenuación versus frecuencia. Con estas dos herramientas, lo primero que hacemos es calcular el presupuesto que va a ser común a los 5 puntos, ya que la relación de potencias es la misma:
\[
  L_{dB} = 10\log_{10}\left(\frac{100}{1}\right)
\]
\[
  L_{dB} = 20\,\mathrm{dB}
\]
Con esto, el segundo paso es analizar cada punto a partir de revisar los gráficos de la Figura 4.3 del libro:
\begin{itemize}
  \item[a)] Para una $f = 300\,\mathrm{kHz}$, podemos ubicar a ojo un valor de $10\,\mathrm{dB}$ de atenuación, con lo cual la longitud $L$ máxima es $L = \text{presupuesto}/\text{tasa de atenuación}$:
  \[
    L = \frac{20\,\mathrm{dB}}{10\,\mathrm{dB/Km}} = 2\,\mathrm{Km}
  \]
  \item[b)] Para $f = 1\,\mathrm{MHz}$, en este caso:
  \[
    L = \frac{20\,\mathrm{dB}}{20\,\mathrm{dB/Km}} = 1\,\mathrm{Km}
  \]
  \item[c)]
  \[
    L_{max} = \frac{20\,\mathrm{dB}}{2,5\,\mathrm{dB/Km}} = 8\,\mathrm{Km}
  \]
  \item[d)]
  \[
    L_{max} = \frac{20\,\mathrm{dB}}{10\,\mathrm{dB/Km}} = 2\,\mathrm{Km}
  \]
  \item[e)] Frecuencia óptima en el caso de la fibra óptica sería el mínimo de la curva, más o menos entre $0,2$ y $0,5$; vamos a usar $0,2$ para el cálculo:
  \[
    L_{max} = \frac{20\,\mathrm{dB}}{0,2\,\mathrm{dB/Km}} = 100\,\mathrm{Km}
  \]
\end{itemize}

  \item
        \textbf{El cable coaxial es un sistema de transmisión con dos conductores. ¿Qué ventaja tiene conectar la malla exterior a tierra?} \\
        Debido al apantallamiento, por construcción, el cable coaxial es mucho menos susceptible que el par trenzado tanto a interferencias como a diafonía. La ventaja de conectar la malla exterior a tierra es que las corrientes inducidas por señales externas se derivan a tierra a través de la malla, reduciendo así el ruido y las interferencias que podrían afectar al conductor interno.

  \item
\textbf{Demuestre que duplicando la frecuencia de transmisión o duplicando la distancia entre las antenas de transmisión y recepción, la potencia recibida se atenúa en 6 dB.} \\
\section*{Demostración}

La pérdida en el espacio libre $L_{dB}$ se define como:
\begin{equation}
    L_{dB} = 20 \log_{10}(f) + 20 \log_{10}(d) - 147,56 \quad [\text{dB}]
 \end{equation}
Donde $f$ es la frecuencia en Hz y $d$ es la distancia en metros.
Alternativamente, analizando la razón de pérdida $L = \frac{P_t}{P_r}$ en escala lineal, tenemos:
\begin{equation}
    L = \left( \frac{4\pi d f}{c} \right)^2
\end{equation}

Expresando esta razón en decibelios ($\text{dB}$):
\begin{equation}
    L_{dB} = 10 \log_{10} \left( \frac{P_t}{P_r} \right) = 20 \log_{10} \left( \frac{4\pi d f}{c} \right)
\end{equation}

A continuación, evaluaremos independientemente cada una de las dos condiciones planteadas en el enunciado.

\subsection*{Caso A: Duplicar la frecuencia ($f' = 2f$)}

Supongamos que mantenemos la distancia $d$ constante y duplicamos la frecuencia de la portadora, de modo que $f' = 2f$. Calculamos la nueva atenuación $L'_{dB}$:

\begin{align*}
    L'_{dB} &= 20 \log_{10} \left( \frac{4\pi d (2f)}{c} \right) \\[0.5em]
    &= 20 \log_{10} \left( 2 \cdot \frac{4\pi d f}{c} \right)
\end{align*}

Aplicando la propiedad de los logaritmos ($\log(a \cdot b) = \log(a) + \log(b)$):

\begin{align*}
    L'_{dB} &= 20 \log_{10}(2) + 20 \log_{10} \left( \frac{4\pi d f}{c} \right) \\[0.5em]
    &= 20 \log_{10}(2) + L_{dB}
\end{align*}

Sabiendo que $20 \log_{10}(2) \approx 20 \cdot 0,30103 = 6,02\text{ dB} \approx 6\text{ dB}$:

\begin{equation}
    L'_{dB} \approx L_{dB} + 6\text{ dB}
\end{equation}

\textbf{Conclusión A:} Al duplicar la frecuencia, la atenuación del medio aumenta en $6\text{ dB}$, lo que equivale a que la potencia recibida $P_r$ se reduce a la cuarta parte, es decir, una pérdida adicional de $6\text{ dB}$.

\subsection*{Caso B: Duplicar la distancia ($d' = 2d$)}

Supongamos ahora que mantenemos la frecuencia $f$ constante y duplicamos la distancia entre antenas, haciendo $d' = 2d$. La nueva atenuación $L''_{dB}$ resulta:

\begin{align*}
    L''_{dB} &= 20 \log_{10} \left( \frac{4\pi (2d) f}{c} \right) \\[0.5em]
    &= 20 \log_{10} \left( 2 \cdot \frac{4\pi d f}{c} \right) \\[0.5em]
    &= 20 \log_{10}(2) + 20 \log_{10} \left( \frac{4\pi d f}{c} \right) \\[0.5em]
    &= 20 \log_{10}(2) + L_{dB}
\end{align*}

De la misma manera:

\begin{equation}
    L''_{dB} \approx L_{dB} + 6\text{ dB}
\end{equation}

\textbf{Conclusión B:} Al duplicar la distancia entre el transmisor y el receptor, la atenuación también se incrementa en $6\text{ dB}$.

\section*{Conclusión}
Queda demostrado que tanto al duplicar la frecuencia de transmisión ($f$) como al duplicar la distancia de enlace ($d$), el término dependiente introduce un factor $20 \log_{10}(2) \approx 6\text{ dB}$ adicional a la pérdida de propagación en espacio libre, atenuando la potencia de la señal recibida en $6\text{ dB}$ en ambos casos. \quad $\blacksquare$
  \item
      Para calcular la longitud típica de la antena, primero se determina la longitud de onda ($\lambda$) de la señal electromagnética a $30\text{ Hz}$:

    \begin{equation}
        c = \lambda \cdot f \implies \lambda = \frac{c}{f}
    \end{equation}

    Sustituyendo la velocidad de la luz $c \approx 3 \times 10^8\text{ m/s}$ y la frecuencia $f = 30\text{ Hz}$:

    \begin{equation}
        \lambda = \frac{3 \times 10^8\text{ m/s}}{30\text{ s}^{-1}} = 10\,000\,000\text{ m} = 10\,000\text{ km}
    \end{equation}

    Dado que la dimensión deseable de la antena debe ser del orden de la mitad de la longitud de onda ($L = \lambda / 2$):

    \begin{equation}
        L = \frac{10\,000\text{ km}}{2} = 5\,000\text{ km}
    \end{equation}

    Por lo tanto, la longitud típica requerida para la antena es de \textbf{$5\,000\text{ km}$} (o $5 \times 10^6\text{ m}$).
  \item 
\textbf{La potencia de la señal de voz está concentrada en torno a los 300 Hz. Las antenas para transmitir esta frecuencia deberían tener un tamaño enormemente grande. Esto hace que, para transmitir voz por radio, la señal deba enviarse modulando una señal de frecuencia superior (portadora) para la que la antena correspondiente requiera un tamaño menor.}

        \textbf{a) ¿Cuál debe ser la longitud de una antena, equivalente a la mitad de la longitud de onda, para enviar una señal de 300 Hz?}

        \textbf{b) Una posible alternativa es emplear algún esquema de modulación, como los descritos en el Capítulo 5, de tal manera que la señal a transmitir tenga un ancho de banda estrecho, centrado en torno a la frecuencia portadora. Supóngase que quisiéramos una antena de 1 metro de longitud. ¿Qué frecuencia de portadora debería utilizarse?}
  
  Para resolver este problema, se utiliza la relación fundamental entre la velocidad de propagación de una onda electromagnética ($c \approx 3 \times 10^8\ \text{m/s}$), la frecuencia ($f$) y la longitud de onda ($\lambda$):

        \[
          \lambda = \frac{c}{f}
        \]

        \textbf{a) Longitud de la antena para una señal de 300 Hz}

        Primero calculamos la longitud de onda:

        \[
          \lambda = \frac{3 \times 10^8\ \text{m/s}}{300\ \text{Hz}}
          = 1\,000\,000\ \text{m} = 1\,000\ \text{km}
        \]

        Como la longitud de la antena debe ser la mitad de la longitud de onda ($L = \lambda/2$):

        \[
          L = \frac{1\,000\,000\ \text{m}}{2}
          = \boxed{500\,000\ \text{m} = 500\ \text{km}}
        \]

        Una antena de 500 km es físicamente irrealizable para un sistema de transmisión práctico, lo que justifica la necesidad de utilizar modulación.

        \medskip
        \textbf{b) Frecuencia portadora requerida para una antena de 1 metro}

        Si la antena mide $L = 1\ \text{m}$ y equivale a la mitad de la longitud de onda ($L = \lambda/2$), la longitud de onda necesaria es:

        \[
          \lambda = 2L = 2 \times 1\ \text{m} = 2\ \text{m}
        \]

        Calculamos la frecuencia portadora:

        \[
          f = \frac{c}{\lambda}
          = \frac{3 \times 10^8\ \text{m/s}}{2\ \text{m}}
          = 150\,000\,000\ \text{Hz}
          = \boxed{150\ \text{MHz}}
        \]

        Para poder utilizar una antena de 1 metro, la señal de voz debe modularse sobre una frecuencia portadora de \textbf{150 MHz}, perteneciente a la banda VHF.

  \item
        Observamos que la longitud de onda completa es $\lambda = 0,005\,\mathrm{m}$, y usando $f = c/\lambda$, despejamos y encontramos que:
        \[
             f = \frac{c}{\lambda} = 6\times10^{10}\,\mathrm{Hz} = 60\,\mathrm{GHz}
        \]
  \item

        Debido a la relación:

        \[
          \frac{P_r}{P_t} = \frac{f^2 A_r A_t}{(cd)^2}
        \]

        Si duplicamos la frecuencia, entonces:

        \[
          \frac{P_r}{P_t} = \frac{(2f)^2 A_r A_t}{(cd)^2}
        \]

        \[
          \frac{P_r}{P_t} = 4 \frac{f^2 A_r A_t}{(cd)^2}; \quad a = \frac{f^2 A_r A_t}{(cd)^2}
        \]

        \[
          \frac{P_r}{P_t} = 4a
        \]

        \[
          P_r = 4a P_t
        \]

        Lo que significa que cuadruplicamos la relación entre potencia recibida y potencia transmitida.

        Por otro lado, si duplicamos el área efectiva de ambas antenas, tenemos que:

        \[
          \frac{P_r}{P_t} = \frac{f^2 \cdot 2A_r \cdot 2A_t}{(cd)^2}
        \]

        O sea:

        \[
          \frac{P_r}{P_t} = 4 \frac{f^2 A_r A_t}{(cd)^2}; \quad a = \frac{f^2 A_r A_t}{(cd)^2}
        \]

        \[
          \frac{P_r}{P_t} = 4a
        \]

        \[
          P_r = 4a P_t
        \]

        Esto demuestra que ambos escenarios llegan a un mismo resultado: una cuadruplicación de la potencia recibida.

  \item \textbf{En la transmisión de radio en el espacio libre, la potencia de la señal se reduce proporcionalmente al cuadrado de la distancia recorrida desde la fuente, mientras que en una transmisión en un cable, la atenuación es una cantidad fija en dB por kilómetro. En la siguiente
tabla se muestra, en dB, la reducción relativa a una referencia dada para la transmisión en
el espacio libre y en un cable uniforme. Rellene las celdas que faltan para completar la
tabla.}

\section*{Deducción de Fórmulas}

\subsection*{1. Transmisión por Radio (Espacio Libre)}
En espacio libre, la atenuación varía con el cuadrado de la distancia ($d^2$). Expresado en escala logarítmica ($\text{dB}$), la atenuación es proporcional a $20 \log_{10}(d)$.

Si se duplica la distancia ($d' = 2d$), el cambio en atenuación es:
\[
\Delta L_{\text{Radio}} = 20 \log_{10}(2) \approx 6,02 \text{ dB} \approx 6 \text{ dB}
\]
Esto significa que **cada vez que la distancia se duplica, la atenuación por radio aumenta en $6 \text{ dB}$ adicionales** (es decir, disminuye $6 \text{ dB}$ respecto a la referencia).

\subsection*{2. Transmisión por Cable}
En un cable, la atenuación es lineal respecto a la distancia ($a \cdot d$). Por lo tanto, en escala en $\text{dB}$, la atenuación es directamente proporcional a la longitud recorrida:
\[
L_{\text{Cable}}(d) = \alpha \cdot d
\]
Dado que para $d = 1\text{ km}$ la atenuación es de $-3\text{ dB}$, la tasa fija de atenuación es de $\alpha = -3\text{ dB/km}$.
Esto implica que **cada vez que la distancia se duplica, la atenuación en $\text{dB}$ también se duplica**.

\section*{Cálculo de Valores}

\begin{itemize}
    \item \textbf{Para $1\text{ km}$:}
    \begin{itemize}
        \item Radio: $-6\text{ dB}$ (Dato de referencia).
        \item Cable: $-3\text{ dB}$ (Dato de referencia).
    \end{itemize}
    
    \item \textbf{Para $2\text{ km}$ (se duplica respecto a $1\text{ km}$):}
    \begin{itemize}
        \item Radio: $-6\text{ dB} - 6\text{ dB} = \mathbf{-12\text{ dB}}$
        \item Cable: $-3\text{ dB} \times 2 = \mathbf{-6\text{ dB}}$
    \end{itemize}

    \item \textbf{Para $4\text{ km}$ (se duplica respecto a $2\text{ km}$):}
    \begin{itemize}
        \item Radio: $-12\text{ dB} - 6\text{ dB} = \mathbf{-18\text{ dB}}$
        \item Cable: $-3\text{ dB} \times 4 = \mathbf{-12\text{ dB}}$
    \end{itemize}

    \item \textbf{Para $8\text{ km}$ (se duplica respecto a $4\text{ km}$):}
    \begin{itemize}
        \item Radio: $-18\text{ dB} - 6\text{ dB} = \mathbf{-24\text{ dB}}$
        \item Cable: $-3\text{ dB} \times 8 = \mathbf{-24\text{ dB}}$
    \end{itemize}

    \item \textbf{Para $16\text{ km}$ (se duplica respecto a $8\text{ km}$):}
    \begin{itemize}
        \item Radio: $-24\text{ dB} - 6\text{ dB} = \mathbf{-30\text{ dB}}$
        \item Cable: $-3\text{ dB} \times 16 = \mathbf{-48\text{ dB}}$
    \end{itemize}
\end{itemize}

\section*{Tabla Completa}

\begin{table}[h!]
\centering
\begin{tabular}{ccc}
\toprule
\textbf{Longitud (km)} & \textbf{Radio (dB)} & \textbf{Cable (dB)} \\
\midrule
1  & $-6$  & $-3$  \\
2  & $-12$ & $-6$  \\
4  & $-18$ & $-12$ \\
8  & $-24$ & $-24$ \\
16 & $-30$ & $-48$ \\
\bottomrule
\end{tabular}

\end{table}

Observamos que para distancias cortas (hasta $4\text{ km}$), el cable atenúa menos que la propagación por radio. A los $8\text{ km}$ ambos medios sufren la misma atenuación ($-24\text{ dB}$), y para distancias mayores ($16\text{ km}$ o más), el medio inalámbrico (radio) resulta ser mucho más eficiente en cuanto a atenuación que el cable. \quad $\blacksquare$

  \item Para la parábola de la Figura 4.12 se quiere demostrar
        que los rayos que parten del foco y se reflejan en la parábola siguen
        trayectorias paralelas al eje \(x\).

        \textbf{a)}

        La ecuación de la parábola es:

        \[
        y^2=2px
        \]

        Despejando \(x\):

        \[
        x=\frac{y^2}{2p}
        \]

        Derivando:

        \[
        \frac{dx}{dy}=\frac{y}{p}
        \]

        Por lo tanto:

        \[
        \frac{dy}{dx}=\frac{p}{y}
        \]

        En el punto \(P(x_1,y_1)\), la pendiente de la tangente \(M\) es:

        \[
        m_M=\frac{p}{y_1}
        \]

        Como la pendiente de una recta es igual a la tangente del ángulo
        que forma con el eje \(x\):

        \[
        \boxed{\tan\beta=\frac{p}{y_1}}
        \]

        \textbf{b)}

        El foco de la parábola se encuentra en:

        \[
        F\left(\frac{p}{2},0\right)
        \]

        La pendiente de la recta \(PF\) es:

        \[
        m_{PF}
        =
        \frac{y_1-0}{x_1-\frac{p}{2}}
        \]

        Como:

        \[
        x_1=\frac{y_1^2}{2p}
        \]

        se obtiene:

        \[
        m_{PF}
        =
        \frac{y_1}
        {\frac{y_1^2}{2p}-\frac{p}{2}}
        \]

        \[
        m_{PF}
        =
        \frac{2py_1}{y_1^2-p^2}
        \]

        \[
        \tan\theta = m_{PF}
        \]

        Aplicando la fórmula de la tangente de la diferencia de dos ángulos:

        \[
        \tan\alpha=
        \frac{\tan\beta-\tan\theta}
        {1+\tan\beta\tan\theta}
        \]

        y utilizando las pendientes de \(M\) y \(PF\), resulta:

        \[
        \tan\alpha=-\frac{p}{y_1}
        \]

        Tomando el valor del ángulo entre las dos rectas:

        \[
        \boxed{\tan\alpha=\frac{p}{y_1}}
        \]

        Como:

        \[
        \tan\alpha=\tan\beta
        \]

        se cumple que:

        \[
        \boxed{\alpha=\beta}
        \]

        Por lo tanto, debido a que el ángulo de incidencia es igual al
        ángulo de reflexión, los rayos reflejados por la parábola siguen
        trayectorias paralelas al eje \(x\).

  \item
\textbf{A menudo es más conveniente expresar las distancias en km en lugar de en m y las frecuencias en MHz en lugar de Hz. Rescriba la Ecuación (4.1) usando estas unidades.}
  
  Para reescribir la fórmula de pérdida en el espacio libre utilizando la frecuencia en megahercios ($f_{\text{MHz}}$) y la distancia en kilómetros ($d_{\text{km}}$), partimos de la ecuación con la frecuencia $f$ expresada en hercios y la distancia $d$ en metros:

        \[
          L_{\text{dB}} = 20\log_{10}(f) + 20\log_{10}(d) - 147{,}56
        \]

        \textbf{Relación de unidades:}
        \[
          f = f_{\text{MHz}} \times 10^6,
          \qquad d = d_{\text{km}} \times 10^3
        \]

        Sustituyendo en la ecuación original:

        \[
          L_{\text{dB}} = 20\log_{10}(f_{\text{MHz}} \times 10^6)
          + 20\log_{10}(d_{\text{km}} \times 10^3) - 147{,}56
        \]

        Aplicando la propiedad $\log(ab) = \log(a) + \log(b)$:

        \[
          \begin{aligned}
            L_{\text{dB}} ={}& 20\log_{10}(f_{\text{MHz}}) + 20\log_{10}(10^6) \\
                            &+ 20\log_{10}(d_{\text{km}}) + 20\log_{10}(10^3) - 147{,}56
          \end{aligned}
        \]

        Calculamos las constantes logarítmicas:

        \[
          20\log_{10}(10^6) = 20 \times 6 = 120,
          \qquad 20\log_{10}(10^3) = 20 \times 3 = 60
        \]

        Por lo tanto:

        \[
          L_{\text{dB}} = 20\log_{10}(f_{\text{MHz}})
          + 20\log_{10}(d_{\text{km}}) + (120 + 60 - 147{,}56)
        \]

        Expresando la frecuencia en MHz y la distancia en km, la pérdida en el espacio libre, en dB, resulta:

        \[
          \boxed{L_{\text{dB}} = 20\log_{10}(f_{\text{MHz}})
          + 20\log_{10}(d_{\text{km}}) + 32{,}44}
        \]

  \item
        \textbf{Suponga que un transmisor emite 50 W de potencia.}

        \textbf{a) Exprese la potencia transmitida en dBm y dBW.}
        $$
          Potencia_{dBW}=10\log\left(\frac{Potencia_W}{1W}\right)
        $$
        $$
          \boxed{
            Potencia_{dBW}=16.99\,dBW
          }
        $$
        $$
          Potencia_{dBm}=10\log\left(\frac{Potencia_{mW}}{1mW}\right)
        $$
        $$
          Potencia_{dBm}=10\log\left(\frac{50000mW}{1mW}\right)
        $$
        $$
          \boxed{Potencia_{dBm}=46.99\,dBm}
        $$

        \textbf{b) Si la potencia del transmisor se aplica a una antena con ganancia unidad, usando una frecuencia de portadora de 900 MHz, ¿cuál es la potencia recibida, en dBm, en el espacio libre a una distancia de 100 m?}

        Para la antena isotrópica ideal, la pérdida en el espacio libre es
        $$
          L_{dB}=20\log(f)+20\log(d)-147.56\,dB
        $$
        $$
          L_{dB}=20\log(900\times10^6)+20\log(100)-147.56
        $$
        $$
          \boxed{L_{dB}=71.52\,dB}
        $$

        Luego la potencia de la señal en la antena receptora, segun la perdida en el espacio libre es
        $$
          P_r=P_T\left(\frac{\lambda}{4\pi d}\right)^2
        $$
        $$
          10\log(P_r)=10\log(P_t) +10\log\left[  \left(\frac{\lambda}{4\pi d}\right)^2\right]
        $$

        Como la pérdida en espacio libre es:
        $$
          L_{dB}=10\log\left(\frac{P_t}{P_r}\right)
        $$

        entonces:
        $$
          P_r(dBm)=P_t(dBm)-L_{dB}
        $$
        $$
          P_r=46.99-71.52
        $$
        $$
          \boxed{P_r=-24.53\,dBm}
        $$

        \textbf{c) Repita el Apartado (b) para una distancia de 10 km.}
        $$
          L_{dB}=20\log(900\times10^6)+20\log(10000)-147.56=111.52\,dB
        $$
        $$
          \boxed{P_r=46.99\,dBm-111.52\,dB=-64.53\,dBm}
        $$

        \textbf{d) Repita (c) pero suponiendo una ganancia para la antena de recepción de 2}

        Para antenas teniendo en cuenta la ganancia, la ecuacion de pérdida en el espacio libre es:
        $$
          \frac{P_t}{P_r}=\frac{(4\pi)^2d^2}{G_rG_t\lambda^2}
        $$

        Podemos rescribir esta ecuacion para obtener la potencia recibida en dBm y obtenemos:
        $$
          P_r(dBm)=P_t(dBm)+G_t(dB)+G_r(dB)-L_{dB}
        $$

        En nuestro caso $G_t=1$ y $G_r=2$. Convertimos esa ganancia a dB:
        $$
          G_r(dB)=10\log(2)
        $$
        $$
          \boxed{G_r=3.01\,dB}
        $$

        Entonces, usando el balance en dB:
        $$
          P_r(dBm)=P_t(dBm)+G_t(dB)+G_r(dB)-L_{dB}
        $$

        Por tanto:
        $$
          P_r=46.99+0+3.01-111.52
        $$
        $$
          \boxed{P_r=-61.52\,dBm}
        $$

  \item
         \begin{itemize}
      \item Potencia de transmisión: $P_t = 0{,}1,\text{W} = 100 \text{mW}$
      \item Frecuencia de operación: $f = 2\text {GHz} = 2 \times 10^9\text{Hz}$
      \item Diámetro de las antenas: $d = 1{,}2 \text{m}$, por lo que $r = 0{,}6,\text{m}$
      \item Distancia del enlace: $d_{\text{dist}} = 24\text {km} = 24,000\text{m}$
      \item Velocidad de propagación: $c \approx 3 \times 10^8\text{m/s}$
      \end{itemize}

      Primero calculamos la longitud de onda:

      \begin{equation}
      \lambda = \frac{c}{f}
      = \frac{3 \times 10^8}{2 \times 10^9}
      = 0{,}15,\text{m}
      \end{equation}

      El área física de la antena es:

      \begin{equation}
      A = \pi r^2
      = \pi(0{,}6)^2
      \approx 1{,}131,\text{m}^2
      \end{equation}

      Considerando una eficiencia de $0{,}56$, el área efectiva resulta:

      \begin{equation}
      A_e = 0{,}56 A
      = 0{,}56 \times 1{,}131
      \approx 0{,}6333,\text{m}^2
      \end{equation}

      \subsection*{a) Ganancia de cada antena}

      La ganancia se obtiene a partir del área efectiva:

      \begin{equation}
      G = \frac{4\pi A_e}{\lambda^2}
      = \frac{4\pi(0{,}6333)}{(0{,}15)^2}
      \approx 353{,}72
      \end{equation}

      Para expresarla en decibelios:

      \begin{equation}
      G_{\text{dB}} = 10\log_{10}(353{,}72)
      \approx 25{,}49,\text{dB}
      \end{equation}

      Por lo tanto, la ganancia de cada antena es aproximadamente $25{,}49,\text{dB}$.

      \subsection*{b) Potencia Efectiva Radiada (EIRP)}

      La EIRP se obtiene multiplicando la potencia transmitida por la ganancia de la antena transmisora:

      \begin{equation}
      \text{EIRP} = P_t G_t
      = 0{,}1 \times 353{,}72
      \approx 35{,}37,\text{W}
      \end{equation}

      La potencia transmitida expresada en dBm es:

      \begin{equation}
      P_{t,\text{dBm}}
      = 10\log_{10}(100)
      = 20,\text{dBm}
      \end{equation}

      Entonces:

      \begin{equation}
      \text{EIRP}{\text{dBm}}
      = P{t,\text{dBm}} + G_{t,\text{dB}}
      = 20 + 25{,}49
      = 45{,}49,\text{dBm}
      \end{equation}

      Por lo tanto, la EIRP es de aproximadamente $35{,}37,\text{W}$, equivalente a $45{,}49,\text{dBm}$.

      \subsection*{c) Potencia recibida ($P_r$)}

      Primero calculamos la pérdida de propagación en el espacio libre:

      \begin{equation}
      L_{\text{dB}}
      = 20\log_{10}\left(
      \frac{4\pi d_{\text{dist}}}{\lambda}
      \right)
      \end{equation}

      Reemplazando:

      \begin{equation}
      L_{\text{dB}}
      = 20\log_{10}\left(
      \frac{4\pi(24,000)}{0{,}15}
      \right)
      \approx 126{,}07,\text{dB}
      \end{equation}

      Finalmente, aplicamos el balance de enlace:

      \begin{equation}
      P_{r,\text{dBm}}
      = P_{t,\text{dBm}}
      + G_{t,\text{dB}}
      + G_{r,\text{dB}}
      - L_{\text{dB}}
      \end{equation}

      Como ambas antenas tienen la misma ganancia:

      \begin{equation}
      P_{r,\text{dBm}}
      = 20 + 25{,}49 + 25{,}49 - 126{,}07
      \approx -55{,}09,\text{dBm}
      \end{equation}

      Por lo tanto, la potencia de la señal a la salida de la antena receptora es aproximadamente $-55{,}09,\text{dBm}$.

  \item
        Para obtener la expresión, comenzamos armando el triángulo e identificando los lados:
        \begin{itemize}
          \item $OB = R$ (radio de la Tierra).
          \item $OA = R + h$ (radio más la altura de la antena).
          \item $AB = d$.
        \end{itemize}
        Planteamos Pitágoras:
        \[
          OA^2 = OB^2 + AB^2
        \]
        \[
          (R+h)^2 = R^2 + d^2
        \]
        Despejamos $d^2$:
        \[
          d^2 = (R+h)^2 - R^2 = R^2 + 2Rh + h^2 - R^2 = 2Rh + h^2
        \]
        Como $h^2$ termina siendo despreciable frente al término $2Rh$ (siendo $R = 6.370$ kilómetros $= 6.370.000$ metros), nos queda:
        \[
          d \approx \sqrt{2Rh}
        \]
        \[
          d(\mathrm{km}) = \sqrt{12,74\cdot h} \approx 3,57\sqrt{h}
        \]
        En el enunciado del ejercicio, vemos que el término $3,57$ está mal escrito, apareciendo como $3,75$; pero en el desarrollo, y previamente en el mismo libro, se menciona que es $3,57$.

  \item

        Para tener en cuenta el peor caso, suponemos que la antena receptora se encuentra a una altura de 0 m.

        La distancia de transmisión será entonces igual a:

        \[
          d = 3,57\ \frac{\text{km}}{\sqrt{\text{m}}} \sqrt{Kh}
        \]

        Donde $K$ es el factor de ajuste que tiene en cuenta la refracción.

        Despejando:

        \[
          h = \left(\frac{d}{3,57\ \frac{\text{km}}{\sqrt{\text{m}}} \sqrt{K}}\right)^2
        \]

        $d = 80\ \text{km}$ y $K \approx \frac{4}{3}$.

        \[
          h = \left(\frac{80\ \text{km}}{3,57\ \frac{\text{km}}{\sqrt{\text{m}}} \sqrt{\frac{4}{3}}}\right)^2
        \]

        \[
          h = 376,63\ \text{m}
        \]

  \item
    El rayo incide sobre la superficie del agua formando un ángulo de $30^\circ$ respecto del horizonte. Como la Ley de Snell utiliza los ángulos medidos respecto de la normal a la superficie, primero calculamos el ángulo de incidencia:

    \begin{equation}
    \theta_1 = 90^\circ - 30^\circ = 60^\circ
    \end{equation}

    Los índices de refracción son:

    \begin{equation}
    n_1 = 1{,}0003
    \end{equation}

    para el aire, y

    \begin{equation}
    n_2 = \frac{4}{3} \approx 1{,}3333
    \end{equation}

    para el agua.

    Aplicando la Ley de Snell:

    \begin{equation}
    n_1 \sin(\theta_1) = n_2 \sin(\theta_2)
    \end{equation}

    Despejamos $\theta_2$:

    \begin{equation}
    \sin(\theta_2) =
    \frac{n_1 \sin(\theta_1)}{n_2}
    \end{equation}

    Reemplazando los valores:

    \begin{equation}
    \sin(\theta_2) =
    \frac{1{,}0003 \sin(60^\circ)}{1{,}3333}
    \approx 0{,}6497
    \end{equation}

    Por lo tanto:

    \begin{equation}
    \theta_2 = \arcsin(0{,}6497) \approx 40{,}52^\circ
    \end{equation}

    Entonces, el rayo se refracta dentro del agua formando un ángulo de aproximadamente $40{,}52^\circ$ respecto de la normal.

    Si se quiere expresar el ángulo respecto de la superficie del agua:

    \begin{equation}
    90^\circ - 40{,}52^\circ = 49{,}48^\circ
    \end{equation}

    Por lo tanto, el ángulo de refracción es de $40{,}52^\circ$ respecto de la normal, o equivalentemente $49{,}48^\circ$ respecto de la superficie del agua.


\end{enumerate}

\end{document}
